% Lösungen Einheit 2 von 3 – Bayes (zusammenbau v0.1)
\einheitenkopf{Einheit 2 von 3 · Bayes}

\erg{10}{a) mit dem Baum \quad b) $\frac{0{,}1 \cdot 0{,}8}{0{,}1 \cdot 0{,}8 + 0{,}9 \cdot 0{,}1} = \frac{0{,}08}{0{,}17} \approx 0{,}471$ \quad c) die Wahrscheinlichkeit, dass ein defektes Gerät aus dem ersten Werk stammt: im Zähler der Pfad „erstes Werk und defekt“, im Nenner die Summe beider Pfade zu „defekt“ \quad d) $\frac{0{,}06}{0{,}06 + 0{,}07} \approx 0{,}462$, also ja, größer als $0{,}4$ \quad e) Ja-Pfade: $\frac{2}{3} \cdot 0{,}3$ (Frage A, abgeschrieben) und $\frac{1}{3} \cdot 0{,}7$ (Frage B, nie abgeschrieben); $\frac{0{,}2}{0{,}2 + \frac{7}{30}} \approx 0{,}462$ \quad f) A und B haben beide den Blauanteil $\frac{1}{3}$, ihre Pfade ergeben zusammen $\frac{2}{3} \cdot \frac{1}{3}$; im Zähler der Pfad über C. Aus dem Wert folgt $\frac{2}{3(r+2)} = \frac{2}{27}$, also $r = 7$}
\erg{11}{a) Fehler: 95\,\% ist die Wahrscheinlichkeit für „positiv, wenn krank“, gefragt ist die umgekehrte Richtung; richtig $\frac{0{,}04 \cdot 0{,}95}{0{,}04 \cdot 0{,}95 + 0{,}96 \cdot 0{,}05} = \frac{0{,}038}{0{,}086} \approx 0{,}442$ \quad b) Eingetreten ist das Ereignis; jeder Weg, auf dem es eintreten kann, gehört zur neuen Gesamtheit, deshalb zählen alle Pfade zu ihm zusammen \quad c) $P(K \cap T) = 0{,}045$, $P(\overline{K} \cap T) = 0{,}095$, $P(K \cap \overline{T}) = 0{,}005$, $P(\overline{K} \cap \overline{T}) = 0{,}855$ \quad d) $\frac{0{,}009}{0{,}009 + 0{,}0891} \approx 0{,}092$; nur etwa jede elfte Frau mit positivem Befund ist erkrankt – der Befund verlangt eine Nachuntersuchung, keine Diagnose}

11c)

\vierfeldertafel{K}{T}{0.045,0.095,0.14,0.005,0.855,0.86,0.05,0.95,1}

